\chapter{Cứu Lúc Lớp Trầm}
\chapterintro{Lớp trầm không hẳn là lớp hỏng}{Chỉ cần một màn chuyển nhịp nhẹ, lớp có thể sống lại mà không cần ồn ào quá mức.}

\section{Cứu tiết sau giờ thể dục}
\begin{tinhhuong}{Tình huống}
Học sinh vừa từ tiết vận động sang. Người thì còn mệt, người thì còn nói chuyện dở, khó vào nề nếp ngay.
\end{tinhhuong}
\begin{goiydungngay}
Cho cả lớp hít vào, thở ra hai nhịp thật nhanh rồi hỏi: \textit{``Ai dùng đúng ba từ để tả trạng thái lớp mình bây giờ?''}
\begin{itemize}
\item Nhận 3 câu trả lời.
\item Cười một nhịp với lớp.
\item Chuyển sang bài bằng câu: \textit{``Rồi, từ bây giờ mình đổi chế độ.''}
\end{itemize}
\end{goiydungngay}
\begin{cauchot}
Đừng chống lại năng lượng cũ của lớp; hãy bẻ lái nó.
\end{cauchot}
\ngancach

\section{Lớp mệt vì bài dài}
\begin{tinhhuong}{Tình huống}
Bài học dài, thông tin nhiều, lớp không quậy nhưng mặt ai cũng đuối. Nếu tiếp tục một mạch, lớp sẽ trượt khỏi bài.
\end{tinhhuong}
\begin{goiydungngay}
Nói: \textit{``Tạm dừng 40 giây. Mỗi em ghi ra một câu: đoạn nào nãy giờ dễ hiểu nhất?''}
\begin{itemize}
\item Không hỏi đoạn khó nhất trước.
\item Chọn vài câu tích cực để đọc lên.
\item Từ đó gom lại năng lượng cho phần sau.
\end{itemize}
\end{goiydungngay}
\begin{cauchot}
Khi lớp mệt, hãy cho các em thấy mình vẫn đang hiểu được điều gì.
\end{cauchot}
\ngancach

\section{Hai phe tranh nhau nói}
\begin{tinhhuong}{Tình huống}
Lớp không trầm mà đang quá rối. Nhiều em cùng muốn chen vào, tiết học bắt đầu loãng vì ai cũng giữ phần mình.
\end{tinhhuong}
\begin{goiydungngay}
Nói ngay: \textit{``Từ bây giờ, ai nói sau phải bắt đầu bằng: Em nối ý bạn ở chỗ nào.''}
\begin{itemize}
\item Biến tranh nhau nói thành nối ý.
\item Em nào không nối được thì chờ lượt sau.
\item Cả lớp thường tự giảm ồn rất nhanh.
\end{itemize}
\end{goiydungngay}
\begin{cauchot}
Muốn cứu lớp ồn, đừng cắt lời hàng loạt; hãy đổi luật chơi của lời nói.
\end{cauchot}
\ngancach

\section{Cứu lớp lúc trời oi hoặc mưa dài}
\begin{tinhhuong}{Tình huống}
Có những buổi lớp trầm không hẳn vì bài, mà vì thời tiết kéo năng lượng xuống đồng loạt.
\end{tinhhuong}
\begin{goiydungngay}
Mở bằng câu: \textit{``Thời tiết hôm nay đang kéo lớp xuống, vậy mình cần một động tác gì để kéo lên?''}
\begin{itemize}
\item Cho lớp nêu 2 đến 3 đề xuất rất ngắn.
\item Chọn một cách đơn giản nhất và làm ngay.
\item Từ đó chuyển sang bài với cảm giác lớp vừa cùng quyết định.
\end{itemize}
\end{goiydungngay}
\begin{cauchot}
Thừa nhận hoàn cảnh của lớp trước khi dạy tiếp thường hiệu quả hơn giả vờ như không có gì.
\end{cauchot}
\ngancach

\section{Một bàn làm mẫu cho cả lớp tỉnh lại}
\begin{tinhhuong}{Tình huống}
Lớp trầm quá rộng, gọi từng cá nhân thì chậm mà gọi cả lớp thì loãng.
\end{tinhhuong}
\begin{goiydungngay}
Chọn một bàn làm mẫu trước: đọc, nói, tóm hoặc trả lời thật nhanh để tạo nhịp cho phần còn lại.
\begin{itemize}
\item Chọn bàn đang tỉnh vừa phải, không cần bàn xuất sắc nhất.
\item Khen một điểm cụ thể sau khi bàn đó làm.
\item Mời bàn khác làm theo cùng khuôn.
\end{itemize}
\end{goiydungngay}
\begin{cauchot}
Khi một nhóm nhỏ bật được nhịp, cả lớp thường dễ đi theo hơn nhiều.
\end{cauchot}
